\documentclass[11pt,a4paper]{article}

\usepackage[utf8]{inputenc}
\usepackage[spanish]{babel}
\usepackage{amsmath, amssymb, amsthm}
\usepackage{mathtools}
\usepackage{geometry}
\usepackage{booktabs}
\usepackage{hyperref}
\usepackage{cite}
\usepackage{float}
\usepackage{xcolor}
\usepackage{enumitem}
\geometry{margin=2.5cm}

\hypersetup{
  colorlinks=true,
  linkcolor=blue!50!black,
  citecolor=blue!50!black,
  urlcolor=blue!50!black
}

\newtheorem{prop}{Proposición}
\newtheorem{lemma}{Lema}
\newtheorem{defn}{Definición}
\theoremstyle{remark}
\newtheorem*{remark}{Observación}

\title{\textbf{Fundamentos teóricos del HybridVacancyDetector:\\
un método multi-señal para cuantificar vacancias\\
en cristales con alta deformación}}
\author{S.~Bergamin\\[2pt]\small Notas internas para discusión con dirección}
\date{\today}

\begin{document}
\maketitle

\begin{abstract}
\noindent
El análisis de Wigner--Seitz (WS) es la herramienta estándar para contar
defectos puntuales en cristales irradiados, pero se sabe que sobreestima el
número de pares Frenkel en el régimen de alta deformación: átomos
transitoriamente desplazados durante una cascada quedan clasificados como
intersticiales, generando vacancias falsas en sus sitios de origen. Las
alternativas reference-free (CNA adaptativo, PTM, análisis de Voronoi) son
robustas al ruido térmico pero introducen sus propias ambigüedades cuando
las señales son débiles. Este documento expone, con derivaciones desde
primeros principios, la base teórica de un método híbrido que combina seis
indicadores locales independientes mediante una función de consenso lineal
ponderada, complementada por dos penalizaciones físicamente motivadas
(tránsito balístico y coherencia de pares Frenkel). El método reproduce el
conteo de WS de forma exacta en el régimen de bajo daño y degrada en forma
controlada hacia un estimador robusto cuando WS comienza a sobrecontar.
\end{abstract}

\tableofcontents
\bigskip

% =============================================================================
\section{Planteo del problema}
% =============================================================================

\subsection{Definición topológica de vacancia}

Sea $\mathcal{L} = \{\mathbf{R}_k\}_{k=1}^{N_\text{ref}}$ el conjunto de
posiciones atómicas de un cristal de referencia (pristino, termalizado a
temperatura $T_\text{ref}$ pero sin daño). Sea $\mathcal{D} =
\{\mathbf{r}_j\}_{j=1}^{N_\text{dmg}}$ el conjunto de posiciones del mismo
sistema tras un evento de daño (cascada PKA, irradiación, deformación
mecánica).

Idealmente, una \emph{vacancia} es un sitio $\mathbf{R}_k \in \mathcal{L}$
que ha quedado desocupado: el átomo que originalmente vibraba en torno a
$\mathbf{R}_k$ ya no se encuentra allí ni en sus inmediaciones. La
definición operativa requiere precisar qué se entiende por
``inmediaciones''.

\subsection{El criterio Wigner--Seitz y su discontinuidad}

El método WS \cite{Stukowski2010,Nordlund2018} asigna cada átomo
$\mathbf{r}_j \in \mathcal{D}$ a su sitio de referencia más cercano:
\begin{equation}
  \pi_\text{WS}(\mathbf{r}_j) \;=\; \arg\min_{k} \;
  \bigl\| \mathbf{r}_j - \mathbf{R}_k \bigr\|_2,
  \label{eq:ws-assign}
\end{equation}
y define la \emph{ocupancia} del sitio $k$ como
$\Omega_k = \bigl| \{ j : \pi_\text{WS}(\mathbf{r}_j) = k \} \bigr|$. El
número de vacancias y de intersticiales se obtiene entonces por simple
conteo:
\begin{equation}
  N_\text{vac}^\text{WS} \;=\; \bigl|\{ k : \Omega_k = 0 \}\bigr|,
  \qquad
  N_\text{int}^\text{WS} \;=\; \sum_{k : \Omega_k \geq 2}\!(\Omega_k - 1).
  \label{eq:ws-count}
\end{equation}

La asignación~(\ref{eq:ws-assign}) es una proyección discontinua: si un
átomo cruza el borde de su celda de Voronoi original (la celda WS), la
ocupancia salta abruptamente de $1$ a $0$ en el sitio de partida y de $1$
a $2$ en el de destino. Para un par primero-vecino a distancia $d_\text{NN}$,
el borde se encuentra a $d_\text{NN}/2$. En BCC con $a_0 = 2.87$~Å (Fe),
$d_\text{NN}/2 \approx 1.24$~Å, una distancia comparable al desplazamiento
RMS térmico de un átomo a temperaturas $T \gtrsim 600$~K, y mucho menor
que el desplazamiento de un átomo balísticamente impulsado durante una
cascada.

\subsection{La sobreestimación de WS en alta deformación}
\label{sec:ws-overest}

Consideremos un átomo $\mathbf{r}_j(t)$ que, durante el pico ballístico
de una cascada, se desplaza una distancia $\Delta r \in
[d_\text{NN}/2,\, d_\text{NN}]$ del centro de su sitio original
$\mathbf{R}_k$, sin abandonar la primera esfera de coordinación.
Físicamente este átomo no representa un defecto permanente: tras la
fase de relajación retornará a $\mathbf{R}_k$ (o a un sitio vecino) y
recombinará con la pretendida vacancia. Pero la
ecuación~(\ref{eq:ws-assign}) lo asigna al sitio vecino $\mathbf{R}_{k'}$,
que pasa a tener $\Omega_{k'} = 2$, mientras $\Omega_k = 0$. El conteo
WS reporta entonces \emph{un par Frenkel falso}.

El problema es bien conocido en la literatura
\cite{Stoller2000,Soneda2003,Granberg2016}: se manifiesta como una
sobreestimación de $N_\text{vac}^\text{WS}$ y $N_\text{int}^\text{WS}$
proporcional a la cantidad de átomos en tránsito balístico al momento de
muestrear el frame. La aplicación más común --- contar Frenkel pairs
en un frame \emph{relajado} tras la cascada --- mitiga el efecto pero
no lo elimina: si la relajación es incompleta, persisten sub-cascadas
en tránsito.

\subsection{Objetivo del método híbrido}

Buscamos un estimador $N_\text{vac}^\text{H}$ que satisfaga:
\begin{enumerate}[label=(\roman*)]
\item \textbf{Reproducibilidad de WS} en régimen de bajo daño:
      $N_\text{vac}^\text{H} \to N_\text{vac}^\text{WS}$ cuando la
      hipótesis de fondo de WS es válida (vacancias bien aisladas,
      sin átomos en tránsito).
\item \textbf{Robustez a ruido térmico}: $N_\text{vac}^\text{H}$ no
      debe crecer con $T$ en un cristal pristino.
\item \textbf{Corrección de la sobreestimación} en alta deformación:
      $N_\text{vac}^\text{H} \leq N_\text{vac}^\text{WS}$ cuando existen
      átomos en tránsito, y el filtrado debe ser físicamente motivado.
\item \textbf{Disponibilidad opcional de la referencia}: si no hay
      frame pristino, el método debe degradar a un estimador
      reference-free sensato.
\end{enumerate}

% =============================================================================
\section{Espacio de candidatos}
% =============================================================================

\subsection{Unión de fuentes}
\label{sec:candidates}

Definimos el conjunto de \emph{candidatos a vacancia} como
\begin{equation}
  \mathcal{C} \;=\; \mathcal{C}_\text{WS} \,\cup\, \mathcal{C}_\text{grid},
\end{equation}
donde
\begin{align}
  \mathcal{C}_\text{WS}   &= \{\mathbf{R}_k \,:\, \Omega_k = 0\},
  \\
  \mathcal{C}_\text{grid} &= \{\mathbf{c}_m \,:\, \mathbf{c}_m
    \text{ es el centro de un cluster del método de grilla FaVaD}\}.
\end{align}
El segundo término permite capturar voids extendidos cuyos centros no
coinciden con sitios de la red de referencia, así como vacancias en
sistemas donde la red de referencia no existe (amorfos, fronteras de
grano, nanopartículas).

\subsection{Construcción de $\mathcal{C}_\text{grid}$}

Sea $\Delta$ el paso de grilla. Para cada punto
$\mathbf{g}_{ijk} = (x_0 + (i+\tfrac{1}{2})\Delta,\,y_0 + (j+\tfrac{1}{2})\Delta,\,z_0 + (k+\tfrac{1}{2})\Delta)$
se calcula la distancia al átomo más cercano,
$d_\text{near}(\mathbf{g}_{ijk}) = \min_{j} \|\mathbf{g}_{ijk} - \mathbf{r}_j\|$,
mediante una linked-cell list con costo $\mathcal{O}(1)$ por consulta.
Los puntos con $d_\text{near} > \tau_\text{grid}$ se agrupan luego
mediante el algoritmo greedy de FaVaD \cite{Toussaint2020}: se ordenan
por $d_\text{near}$ decreciente, se siembra un cluster nuevo desde el
punto con mayor $d_\text{near}$ disponible, y se absorben todos los
puntos no asignados a distancia $\leq r_\text{cluster}$.

\paragraph{Elección de $\tau_\text{grid}$.} El argumento físico es
crucial: el sitio octaédrico de una FCC se encuentra a
$d_\text{oct} = d_\text{NN}/\sqrt{2}\,\approx\,0.71\,d_\text{NN}$ del
átomo más cercano. El sitio tetraédrico está a
$d_\text{tet} = d_\text{NN}\sqrt{3}/2 \approx 0.87\,d_\text{NN}$. Si
$\tau_\text{grid} < d_\text{oct}$, la grilla detectará cada hueco
intersticial de la red pristina como vacancia espuria, generando
$\mathcal{O}(N_\text{ref})$ falsos positivos. La elección
$\tau_\text{grid} = 0.9\,d_\text{NN}$ asegura que solo voids reales
(con al menos un sitio atómico faltante) disparen la detección.

\subsection{Deduplicación geométrica}

Si $\mathbf{c}_m \in \mathcal{C}_\text{grid}$ se encuentra a distancia
menor que $r_\text{dedup} = d_\text{NN}/2$ de algún
$\mathbf{R}_k \in \mathcal{C}_\text{WS}$, se descarta $\mathbf{c}_m$
y se conserva $\mathbf{R}_k$ (con su \texttt{ref\_site\_idx}
asociado). La búsqueda se realiza con una segunda linked-cell list
sobre $\mathcal{C}_\text{WS}$ para costo total
$\mathcal{O}(|\mathcal{C}_\text{grid}|)$.

% =============================================================================
\section{Las seis señales positivas}
% =============================================================================

Para cada candidato $\mathbf{x} \in \mathcal{C}$ se computan seis
indicadores escalares $s_\alpha(\mathbf{x}) \in [0,1]$, donde
$s_\alpha \to 1$ indica fuerte evidencia de vacancia. Cada señal
está respaldada por un argumento físico independiente.

% -----------------------------------------------------------------------------
\subsection{Señal 1: ocupancia WS binaria}
\label{sec:sig-ws}

\begin{equation}
  s_\text{WS}(\mathbf{x}) \;=\;
  \begin{cases}
    1 & \text{si } \mathbf{x} = \mathbf{R}_k \text{ y } \Omega_k = 0,\\
    0 & \text{en caso contrario}.
  \end{cases}
  \label{eq:sig-ws}
\end{equation}

\textbf{Justificación.} Captura literalmente la decisión WS para
preservar el invariante (i) de la sección~\ref{sec:ws-overest}.
\textbf{Complejidad.} $\mathcal{O}(1)$ por candidato (lookup).

% -----------------------------------------------------------------------------
\subsection{Señal 2: ocupancia gaussiana suave (Soft-WS)}
\label{sec:sig-softws}

La ocupancia binaria de WS sufre la discontinuidad descrita en
\S\ref{sec:ws-overest}. La definimos en versión suave como
\begin{equation}
  \Omega_\text{soft}(\mathbf{x}) \;=\;
  \frac{1}{Z(\sigma_T)} \sum_{j=1}^{N_\text{dmg}}
  \exp\!\left( -\frac{\|\mathbf{x} - \mathbf{r}_j\|^2}{2\sigma_T^2} \right),
  \qquad
  s_\text{soft}(\mathbf{x}) \;=\; \max\!\bigl(0,\, 1 - \Omega_\text{soft}(\mathbf{x})\bigr).
  \label{eq:sig-soft}
\end{equation}

\paragraph{Derivación del factor de normalización $Z$.}
En un cristal pristino infinito de densidad $\rho_0 = N/V$, la ocupancia
esperada en cualquier punto $\mathbf{x}$ debe ser $1$. Modelamos la
distribución espacial de átomos como un campo continuo de densidad
$\rho_0$ y calculamos:
\begin{equation}
  \mathbb{E}\!\left[\sum_j e^{-r_j^2/2\sigma^2}\right]
  \;=\; \rho_0 \int_{\mathbb{R}^3} e^{-r^2/2\sigma^2}\,\mathrm{d}^3 r
  \;=\; \rho_0\,(2\pi\sigma^2)^{3/2}.
\end{equation}
Imponiendo $\Omega_\text{soft}^\text{pristino} = 1$ se obtiene
\begin{equation}
  Z(\sigma_T) \;=\; \rho_0\,(2\pi\sigma_T^2)^{3/2}.
  \label{eq:sig-soft-Z}
\end{equation}

\paragraph{Límites asintóticos.}
\begin{itemize}
\item $\sigma_T \to 0$: el peso gaussiano colapsa a una delta y $\Omega_\text{soft}$
      cuenta exactamente $1$ si hay un átomo en $\mathbf{x}$, $0$ en caso
      contrario. Recupera la decisión WS binaria.
\item $\sigma_T \gg d_\text{NN}$: $\Omega_\text{soft}$ tiende a una constante
      $\sim 1$ y la señal pierde toda discriminación. Esto fija el
      régimen útil: $\sigma_T \sim 0.05\text{--}0.20\,d_\text{NN}$.
\end{itemize}

\paragraph{Conexión con la estadística clásica.} La fórmula
(\ref{eq:sig-soft})--(\ref{eq:sig-soft-Z}) es una estimación de densidad
de tipo Parzen--Rosenblatt \cite{Silverman1986}: usamos un kernel
gaussiano con ancho de banda $h = \sigma_T$ para inferir la
densidad atómica local. Si el campo $\rho(\mathbf{x})$ está bien
estimado, $\rho(\mathbf{x})/\rho_0$ es una probabilidad de ocupación
generalizada. Esta es la relación formal entre la señal Soft-WS y las
señales de densidad usadas más adelante (\S\ref{sec:sig-density}).

\textbf{Complejidad.} $\mathcal{O}(N_{3\sigma})$ por candidato, donde
$N_{3\sigma}$ es el número de vecinos a distancia $\leq 3\sigma_T$.

% -----------------------------------------------------------------------------
\subsection{Señal 3: anomalía de primer vecino (proxy de Voronoi)}
\label{sec:sig-voronoi}

\paragraph{Motivación.} La descomposición de Voronoi
\cite{Stukowski2010,Brostow1978} asigna a cada átomo un poliedro
$\mathcal{V}_j$ que contiene todos los puntos del espacio más
cercanos a $\mathbf{r}_j$ que a cualquier otro átomo. En presencia de
una vacancia, los poliedros vecinos se dilatan: el volumen
$V(\mathcal{V}_j)$ excede su valor de referencia. La señal de Voronoi
es entonces
\begin{equation}
  s_\text{vor}^\text{ideal}(\mathbf{x}) \;=\; \mathrm{clip}_{[0,1]}\!\left(
  \frac{\bar V_\text{loc}(\mathbf{x}) - V_0}{V_0} \right),
\end{equation}
donde $V_0$ es el volumen Voronoi de referencia y $\bar V_\text{loc}$
es el promedio sobre los átomos en la primera esfera de $\mathbf{x}$.

\paragraph{Proxy implementado.} Calcular $V(\mathcal{V}_j)$ exigiría
incorporar una librería de Voronoi tridimensional (e.g.\ \texttt{voro++}
\cite{Rycroft2009}). Usamos en su lugar un \emph{proxy del primer
momento}: para cada átomo $\mathbf{r}_j$ a distancia $\leq 2\,d_\text{NN}$
de $\mathbf{x}$, registramos su propia distancia al vecino más cercano
$d_{\text{NN},j}^\text{local}$. La señal es
\begin{equation}
  s_\text{vor}(\mathbf{x}) \;=\; \mathrm{clip}_{[0,1]}\!\left(
  \frac{\mathrm{med}_j \,d_{\text{NN},j}^\text{local}
        \;-\; d_\text{NN}}{d_\text{NN}} \right).
  \label{eq:sig-vor}
\end{equation}

\paragraph{Justificación del proxy.} El radio inscripto de
$\mathcal{V}_j$ está acotado superiormente por
$d_{\text{NN},j}^\text{local}/2$, y para una red regular la mediana de
$d_{\text{NN},j}^\text{local}$ coincide con $d_\text{NN}$. En presencia
de una vacancia, las distancias $d_{\text{NN},j}^\text{local}$ de los
átomos del primer shell se dilatan en proporción a la relajación
elástica de la vecindad. La proporcionalidad
$(d_\text{NN}^\text{local} - d_\text{NN})/d_\text{NN}
\sim (V_\text{loc} - V_0)/(3 V_0)$ se obtiene a primer orden por la
relación volumen--longitud $V \propto d^3 \Rightarrow \delta V/V =
3\,\delta d/d$.

\textbf{Limitación.} El proxy captura el primer momento de la
distribución, no los efectos de asimetría. Para sistemas con
distorsiones tetragonales o trigonales fuertes la señal puede
infraestimar la anomalía; en metales BCC/FCC sin texturas la captura
es buena.

% -----------------------------------------------------------------------------
\subsection{Señal 4: déficit de densidad local (KDE)}
\label{sec:sig-density}

\begin{equation}
  \rho(\mathbf{x}) \;=\;
  \frac{1}{(2\pi h^2)^{3/2}} \sum_{j: \|\mathbf{x}-\mathbf{r}_j\| < 3h}
  \exp\!\left( -\frac{\|\mathbf{x} - \mathbf{r}_j\|^2}{2 h^2} \right),
  \qquad
  s_\text{dens}(\mathbf{x}) \;=\; \mathrm{clip}_{[0,1]}\!\left(
  1 - \frac{\rho(\mathbf{x})}{\rho_0}\right).
  \label{eq:sig-density}
\end{equation}

\paragraph{Selección del ancho de banda $h$.} La regla de Silverman
para densidad multivariada en $d$ dimensiones \cite{Silverman1986} da
$h_\text{opt} \sim n^{-1/(d+4)} \sigma$. Para $d=3$ y
$\sigma \approx d_\text{NN}/2$ se obtiene un ancho de banda óptimo del
orden de $d_\text{NN}/4$. Adoptamos
$h = \max(d_\text{NN}/4,\,0.5\,\text{Å})$ para robustez en cristales
con $d_\text{NN}$ pequeño.

\paragraph{Comparación con Soft-WS.} Las señales \ref{eq:sig-soft} y
\ref{eq:sig-density} son matemáticamente análogas: ambas son
estimadores de densidad de tipo Parzen, con anchos de banda distintos
($\sigma_T$ vs $h$). La diferencia operativa es la escala física:
$\sigma_T$ captura la fluctuación térmica de un sitio individual,
$h$ captura el campo de densidad promedio. Su redundancia controlada
es deseable: en condiciones donde ambos coinciden el consenso es
robusto; cuando discrepan (vibración térmica anómala localizada vs
vacío extendido) la combinación permite distinguir el origen.

% -----------------------------------------------------------------------------
\subsection{Señal 5: anomalía SOAP del entorno}
\label{sec:sig-soap}

Sea $\tilde q^j$ el descriptor SOAP normalizado del átomo $j$ y
$\bar q(T)$ el descriptor medio del cristal pristino a la misma
temperatura. Se define la distancia FaVaD
$d^j = \|\tilde q^j - \bar q(T)\|_2$, que en un cristal pristino sigue
una distribución $\chi$ con $k$ grados de libertad
\cite{Toussaint2020}:
\begin{equation}
  P(d \,|\, k,\sigma) \;\propto\; d^{k-2}\,
  \exp\!\left(-\frac{d^2}{2\sigma^2}\right).
\end{equation}
La señal SOAP del candidato $\mathbf{x}$ es el promedio normalizado
sobre los vecinos del primer shell:
\begin{equation}
  \bar d_\text{shell}(\mathbf{x}) \;=\;
  \frac{1}{|\mathcal{N}(\mathbf{x})|}
  \sum_{j \in \mathcal{N}(\mathbf{x})} d^j,
  \qquad
  s_\text{SOAP}(\mathbf{x}) \;=\; \mathrm{clip}_{[0,1]}\!\left(
  \frac{\bar d_\text{shell}(\mathbf{x})}{\langle d\rangle + 3\sqrt{\mathrm{Var}(d)}}\right),
  \label{eq:sig-soap}
\end{equation}
donde $\mathcal{N}(\mathbf{x}) = \{j : \|\mathbf{x} - \mathbf{r}_j\| <
1.2\,d_\text{NN}\}$ y $\langle d \rangle$, $\mathrm{Var}(d)$ son la
media y varianza muestrales de $d$ en el frame de referencia.

\paragraph{Interpretación estadística.} La normalización por
$\langle d \rangle + 3\sqrt{\mathrm{Var}(d)}$ corresponde a un
$z$-score de tres sigmas: la señal alcanza el valor $1$ cuando la
distancia promedio de los vecinos supera tres desviaciones estándar
del modo de la distribución $\chi$, lo cual ocurre con probabilidad
$\sim 10^{-3}$ en un cristal pristino. Es decir, $s_\text{SOAP} \to 1$
solamente cuando la evidencia estadística de distorsión es fuerte.

\textbf{Sensibilidad.} Esta es la única señal que aprovecha la
estructura química local más allá de las distancias geométricas, y
es la más informativa en cristales con orden orientacional complejo
(HEA, intermetálicos, fases de Laves).

% -----------------------------------------------------------------------------
\subsection{Señal 6: anomalía topológica de coordinación}
\label{sec:sig-topology}

Sea $n_j = |\{i : \|\mathbf{r}_i - \mathbf{r}_j\| < 1.2\,d_\text{NN},
\, i\neq j\}|$ el número de primeros vecinos del átomo $j$. Para una
red ideal sin defectos, $n_j$ es una constante $n_\text{ref}$ ($=8$
para BCC, $12$ para FCC, $12$ para HCP en su primera capa). Se define
\begin{equation}
  s_\text{topo}(\mathbf{x}) \;=\; \mathrm{clip}_{[0,1]}\!\left(
  \frac{1}{|\mathcal{N}(\mathbf{x})|}
  \sum_{j \in \mathcal{N}(\mathbf{x})}
  \frac{\bigl|n_j - n_\text{ref}\bigr|}{n_\text{ref}}\right).
  \label{eq:sig-topo}
\end{equation}

\paragraph{Conexión con CNA y PTM.} Las técnicas modernas de
identificación estructural --- Adaptive Common Neighbor Analysis
(a-CNA) \cite{Stukowski2012} y Polyhedral Template Matching
(PTM) \cite{Larsen2016} --- usan caracterizaciones más sofisticadas
del entorno local: a-CNA computa la firma topológica
$(j_1, j_2, j_3)$ de cada par vecino; PTM ajusta el polígono local a
una plantilla ideal mediante minimización de RMSD. Nuestra señal
$s_\text{topo}$ es una versión escalarizada y simplificada del
número de coordinación, que sacrifica la discriminación BCC/FCC/HCP
para ganar simplicidad y velocidad. En BCC puro o FCC puro la
discriminación es innecesaria, y en aleaciones HEA la estructura
nominal está fijada por la fase mayoritaria.

\paragraph{Sensibilidad a una vacancia simple.} Una vacancia en BCC
elimina uno de los ocho primeros vecinos de cada átomo del shell,
dando $|n_j - 8|/8 = 1/8 = 0.125$. En FCC el efecto es
$1/12 = 0.083$. La señal es por lo tanto débil para vacancias
aisladas en redes de alta coordinación y fuerte en redes de baja
coordinación o en presencia de voids extendidos.

% =============================================================================
\section{Las dos penalties}
% =============================================================================

Mientras las seis señales positivas suman evidencia a favor de
clasificar un candidato como vacancia, las penalties restan
evidencia cuando hay indicios de que la vacancia es transitoria o
parte de un par Frenkel próximo a recombinar.

% -----------------------------------------------------------------------------
\subsection{Penalty 1: filtro de tránsito balístico}
\label{sec:pen-transit}

\paragraph{Motivación cinemática.} Tras la fase ballística de una
cascada, un átomo desplazado tiene tres destinos posibles:
(a) retornar a su sitio original (recombinación);
(b) anclarse en una posición intersticial estable o cuasi-estable;
(c) quedar atrapado por un sumidero (frontera de grano, dislocación).
En los casos (a) y (b), durante un intervalo $\tau_\text{rel}$
posterior al pico térmico, el átomo se encuentra a distancia
$\Delta r \in [0.5\,d_\text{NN},\,1.5\,d_\text{NN}]$ de su sitio
nominal: lo suficientemente lejos para haber abandonado su celda WS,
lo suficientemente cerca para no haberse anclado en otro sitio
intersticial estable. Una vacancia reportada por WS cuyo vecino más
próximo sea un átomo en este rango es candidato a ser un falso
positivo.

\paragraph{Definición operativa.} Dado un candidato $\mathbf{x}$,
sea $r_\text{min}(\mathbf{x})$ la distancia al átomo más cercano que
satisface dos condiciones simultáneas:
\begin{enumerate}[label=(\arabic*)]
\item $r \in [r_\text{lo}, r_\text{hi}]$ con $r_\text{lo} = 0.5\,d_\text{NN}$
      y $r_\text{hi} = f_\text{T}\cdot d_\text{NN}$ (factor
      $f_\text{T}$ típicamente $1.5$);
\item el átomo está clasificado por WS como \emph{intersticial}
      (i.e.\ pertenece a un sitio con $\Omega \geq 2$).
\end{enumerate}
La penalización es
\begin{equation}
  p_\text{transit}(\mathbf{x}) \;=\;
  \begin{cases}
    \dfrac{r_\text{hi} - r_\text{min}(\mathbf{x})}{r_\text{hi} - r_\text{lo}}
      & \text{si } r_\text{min}(\mathbf{x}) \in [r_\text{lo}, r_\text{hi}],\\
    0 & \text{en caso contrario}.
  \end{cases}
  \label{eq:pen-transit}
\end{equation}
La forma lineal en $r$ es la elección más simple que satisface
$p_\text{transit} \to 1$ cuando el átomo está justo en el borde de la
celda WS original ($r \to r_\text{lo}$) y $p_\text{transit} \to 0$
cuando el átomo está demasiado lejos para ser un tránsito directo
($r \to r_\text{hi}$).

\paragraph{Por qué exigir condición (2).} En un cristal con vacancias
reales, los átomos del primer shell del sitio vacante se encuentran
sistemáticamente a distancias en $[r_\text{lo}, r_\text{hi}]$. Si la
condición (2) se relajara a un simple criterio de distancia o a
``átomo distorsionado'' (\texttt{defect\_prob} > umbral), la penalty
se dispararía para toda vacancia real y se cancelarían todas las
detecciones. La restricción a átomos clasificados WS como
intersticiales es lo que distingue un par Frenkel transitorio de una
vacancia aislada permanente.

\subsection{Penalty 2: coherencia de pares Frenkel}
\label{sec:pen-frenkel}

\paragraph{Motivación termodinámica.} El radio de recombinación
$r_\text{rec}$ de un par Frenkel es la distancia bajo la cual
vacancia e intersticial se aniquilan espontáneamente. Para Fe BCC,
estudios atomísticos lo sitúan en $r_\text{rec} \approx 3.3$~Å
\cite{Bjorkas2009,Becquart2010,Stoller2000}. Para FCC (Cu, Ni, FeCrNi
austenítico) los valores reportados están en $3.0\text{--}3.5$~Å.

\paragraph{Definición.} Sea $r_\text{int}(\mathbf{x})$ la distancia
mínima entre el candidato $\mathbf{x}$ y cualquier sitio WS
clasificado como intersticial. La penalty se modela como un
\emph{cluster Gaussian}:
\begin{equation}
  p_\text{Frenkel}(\mathbf{x}) \;=\;
  \exp\!\left( -\frac{r_\text{int}^2(\mathbf{x})}{2\,r_\text{rec}^2}\right).
  \label{eq:pen-frenkel}
\end{equation}

\paragraph{Comportamiento asintótico.}
\begin{itemize}
\item $r_\text{int} \to 0$: par Frenkel coincidente; $p_\text{Frenkel} = 1$.
\item $r_\text{int} = r_\text{rec}$: $p_\text{Frenkel} = e^{-1/2} \approx 0.61$
      (penalty fuerte pero no descarta automáticamente).
\item $r_\text{int} \gg r_\text{rec}$: $p_\text{Frenkel} \to 0$
      (vacancia aislada, no se penaliza).
\end{itemize}

\paragraph{Diferencia con transit.} La penalty de tránsito (\S\ref{sec:pen-transit})
es un filtro \emph{duro} (binario, alto-bajo) que descarta la vacancia
si dispara por encima de $0.5$. La penalty de Frenkel es \emph{suave}:
resta peso al consenso pero no descarta automáticamente. La razón
física: un par Frenkel a $r > r_\text{rec}$ es estable
termodinámicamente; la vacancia es real. Solamente cuando vacancia e
intersticial están muy próximos ($r < r_\text{rec}/\sqrt{2}$) la
evidencia para su carácter transitorio es decisiva.

% =============================================================================
\section{Función de consenso}
% =============================================================================

\subsection{Definición}
\label{sec:consensus}

Asignamos pesos $w_\alpha \geq 0$ a cada señal positiva
$\alpha \in \{\text{WS}, \text{soft}, \text{vor}, \text{dens},
\text{SOAP}, \text{topo}\}$ y a cada penalty
$\beta \in \{\text{transit}, \text{Frenkel}\}$. El \emph{score de
consenso} es
\begin{equation}
  S(\mathbf{x}) \;=\;
  \frac{1}{\sum_\alpha w_\alpha}
  \sum_\alpha w_\alpha\, s_\alpha(\mathbf{x})
  \;-\; w_\text{transit}\,p_\text{transit}(\mathbf{x})
  \;-\; w_\text{Frenkel}\,p_\text{Frenkel}(\mathbf{x}).
  \label{eq:consensus}
\end{equation}

\subsection{Regla de aceptación}

Un candidato se acepta como vacancia si y solo si:
\begin{equation}
  \boxed{
  \;\text{accept}(\mathbf{x}) \;\Longleftrightarrow\;
  \bigl[s_\text{WS}(\mathbf{x}) \geq \tfrac{1}{2}
        \;\;\vee\;\; S(\mathbf{x}) \geq \tau\bigr]
  \;\wedge\; p_\text{transit}(\mathbf{x}) < \tfrac{1}{2}
  \;\wedge\; p_\text{Frenkel}(\mathbf{x}) < \tfrac{1}{2}.\;}
  \label{eq:accept}
\end{equation}
La cláusula disyuntiva en corchetes admite dos rutas independientes
de aceptación: (a) la vacancia es detectada por WS, o (b) las señales
positivas combinadas alcanzan el umbral $\tau$ sin necesidad de WS.
Las dos cláusulas de penalty descartan candidatos con evidencia
fuerte de tránsito o recombinación, respectivamente.

\subsection{Invariantes demostrables}

\begin{prop}[Reproducibilidad de WS]
\label{prop:ws-invariant}
Si los pesos satisfacen $w_\text{transit} = w_\text{Frenkel} = 0$
y $w_\text{WS} > 0$, $w_\alpha = 0$ para $\alpha \neq \text{WS}$,
entonces
$N_\text{vac}^\text{H} = N_\text{vac}^\text{WS}$ y los conjuntos
de sitios reportados coinciden exactamente.
\end{prop}

\begin{proof}
Con la elección de pesos, $S(\mathbf{x}) = s_\text{WS}(\mathbf{x})$.
La regla (\ref{eq:accept}) se reduce a $s_\text{WS} \geq 1/2$, que
es verdadero sii $\Omega_k = 0$ por la definición (\ref{eq:sig-ws}).
Las cláusulas de penalty son vacuamente verdaderas pues los pesos
están en cero. El conjunto aceptado coincide entonces con
$\{\mathbf{R}_k : \Omega_k = 0\} = \mathcal{C}_\text{WS}$.
\qed
\end{proof}

\begin{prop}[Robustez a ruido térmico]
\label{prop:thermal}
Sea un cristal pristino sometido a desplazamientos térmicos
gaussianos $\mathcal{N}(0, \sigma_T^2 \mathbb{I})$ con
$\sigma_T < d_\text{NN}/4$. Entonces, con probabilidad
$\geq 1 - \mathcal{O}(e^{-d_\text{NN}^2/8\sigma_T^2})$, se tiene
$N_\text{vac}^\text{H} = 0$ para cualquier preset con
$w_\text{WS} > 0$ o $w_\text{soft} > 0$.
\end{prop}

\begin{proof}[Esquema]
Bajo la hipótesis $\sigma_T < d_\text{NN}/4$, ningún átomo cruza el
borde de su celda WS salvo con probabilidad exponencialmente pequeña,
luego $s_\text{WS}(\mathbf{R}_k) = 0$ para todo $k$. Análogamente, la
suma de gaussianas que define $\Omega_\text{soft}$ permanece cercana a
$1$ (con desviación $\mathcal{O}(\sigma_T/d_\text{NN})$), de modo que
$s_\text{soft} \to 0$. Los demás $s_\alpha$ son consistentemente
pequeños por argumentos análogos. El score consensuado no alcanza el
umbral $\tau \geq 0.35$ en ningún candidato. Como
$\mathcal{C} = \emptyset$ con alta probabilidad (no hay sitios WS
vacantes ni clusters de grilla por encima del umbral
$\tau_\text{grid}$), $N_\text{vac}^\text{H} = 0$.
\qed
\end{proof}

\begin{prop}[Filtrado de tránsito balístico]
\label{prop:transit}
Si todos los pares Frenkel reportados por WS son ballísticos (es
decir, el átomo intersticial se encuentra a distancia
$\leq r_\text{hi}$ del sitio vacante asociado), entonces el preset
\emph{robusto} ($w_\text{transit} = 1.5 > 0$) descarta todos los
pares y reporta $N_\text{vac}^\text{H} = 0 < N_\text{vac}^\text{WS}$.
\end{prop}

\begin{proof}
Para cada $\mathbf{R}_k \in \mathcal{C}_\text{WS}$, existe por
hipótesis un átomo intersticial WS a distancia
$\leq r_\text{hi} = f_\text{T}\,d_\text{NN}$. Por la definición
(\ref{eq:pen-transit}), $p_\text{transit}(\mathbf{R}_k) \geq 0$;
la condición $r \in [r_\text{lo}, r_\text{hi}]$ con $r_\text{lo} < r$
implica $p_\text{transit}(\mathbf{R}_k) > 0$; en particular, para
$r < (r_\text{lo} + r_\text{hi})/2$ (mitad del rango) se tiene
$p_\text{transit} > 1/2$. La cláusula
$p_\text{transit} < 1/2$ de (\ref{eq:accept}) descarta el candidato.
\qed
\end{proof}

% =============================================================================
\section{Calibración de parámetros físicos}
% =============================================================================

\subsection{Desplazamiento térmico $\sigma_T$ y la relación de Debye--Waller}

El factor de Debye--Waller \cite{Kittel2004} expresa el desplazamiento
RMS de un átomo en torno a su sitio nominal a temperatura $T$:
\begin{equation}
  \langle u^2 \rangle(T) \;=\;
  \frac{3\hbar^2}{m\,k_\text{B}\Theta_\text{D}}\,
  \left[ \frac{T}{\Theta_\text{D}}\,\phi\!\left(\frac{\Theta_\text{D}}{T}\right)
         + \frac{1}{4}\right],
  \label{eq:debye-waller}
\end{equation}
donde $\Theta_\text{D}$ es la temperatura de Debye, $m$ la masa
atómica, y $\phi(x) = (1/x)\int_0^x \xi/(e^\xi - 1)\,\mathrm{d}\xi$
la función de Debye. En el régimen clásico $T \gg \Theta_\text{D}/2$
la fórmula se reduce a
\begin{equation}
  \langle u^2 \rangle(T) \;\approx\;
  \frac{3\,k_\text{B} T}{m\,\omega_\text{D}^2},
  \qquad \sigma_T = \sqrt{\langle u^2\rangle/3}.
\end{equation}
Para Fe ($\Theta_\text{D} = 470$~K, $m = 56$~amu) a $T = 300$~K:
$\sigma_T \approx 0.07$~Å, comparable al default
$0.05\,d_\text{NN} = 0.12$~Å. Para FeCrNi a $T = 300$~K, valores
experimentales reportados están entre $0.07$ y $0.10$~Å.

\paragraph{Estimación operativa.} En ausencia de un cálculo
Debye--Waller explícito, adoptamos $\sigma_T = 0.05\,d_\text{NN}$ como
default conservador, con flag CLI para override del usuario en casos
de temperatura elevada (donde
$\sigma_T \gtrsim 0.10\,d_\text{NN}$).

\subsection{Radio de recombinación $r_\text{rec}$ y la teoría NRT}

El estándar Norgett--Robinson--Torrens (NRT) \cite{Norgett1975} para
producción de defectos por irradiación supone una distribución
espacial de pares Frenkel con radio de recombinación efectivo
$r_\text{rec}$. Para Fe BCC los cálculos más detallados
\cite{Bjorkas2009,Becquart2010} convergen a
$r_\text{rec} \approx 3.3$~Å, aproximadamente $1.3\,d_\text{NN}$.
Adoptamos este valor como default y exponemos el flag
\texttt{--hybrid-recomb-radius} para ajustar a otros materiales.

\subsection{Factor de tránsito $f_\text{T}$}

La elección $f_\text{T} = 1.5$ corresponde al borde de la primera
zona de cuasi-estabilidad para un átomo en transito balístico: a
distancias mayores, el átomo se encuentra en la cuenca de atracción
del sitio nominal vecino y deja de ser un tránsito directo.
Argumentos análogos a los de
Becquart \& Domain \cite{Becquart2010} sustentan el valor.

% =============================================================================
\section{Comparación con la literatura}
% =============================================================================

\subsection{Análisis de Voronoi (OVITO)}

El módulo Voronoi de OVITO \cite{Stukowski2010Ovito} computa el
volumen completo de cada poliedro $V(\mathcal{V}_j)$ y permite
caracterizar la firma topológica $(f_3, f_4, f_5, f_6, \ldots)$ del
mismo. Es más informativo que nuestra señal $s_\text{vor}$ (que solo
captura el primer momento), pero su costo computacional es
$\mathcal{O}(N \log N)$ con prefactor grande, y la firma topológica
no es directamente útil para contar vacancias (los volúmenes
inflados por una vacancia son ambiguos sin un umbral material).

\subsection{Adaptive Common Neighbor Analysis (a-CNA)}

a-CNA \cite{Stukowski2012} clasifica cada átomo por la firma
topológica $(j_1, j_2, j_3)$ de sus pares vecinos, usando un radio de
corte adaptativo derivado de la distribución local de distancias.
Es el estándar de facto para identificar fases (BCC, FCC, HCP, ico) y
detectar atomos defecto-adyacentes. Nuestra señal $s_\text{topo}$ es
una simplificación: solo cuenta vecinos, no clasifica firmas. La
elección se justifica porque (i) en una aleación HEA mono-fase la
discriminación BCC/FCC ya está fija, y (ii) un detector específico
para vacancias no necesita resolver fases.

\subsection{Polyhedral Template Matching (PTM)}

PTM \cite{Larsen2016} es la técnica más robusta al ruido térmico:
ajusta el polígono local de cada átomo a una plantilla ideal
mediante minimización de RMSD invariante por rotación. Reporta a
qué fase pertenece cada átomo y un score $\delta$ de calidad del
ajuste. Para detección de vacancias se puede usar PTM como
pre-filtro de átomos defecto-adyacentes, y luego un segundo paso de
localización geométrica. Nuestro método es más simple, pero podría
incorporar PTM como una sexta señal alternativa a $s_\text{topo}$
sin cambios en la arquitectura.

\subsection{Wigner--Seitz adaptativo y promedios temporales}

Una alternativa al método aquí presentado es \emph{relajar la malla
WS}: en lugar de usar las posiciones $\mathbf{R}_k$ fijas, se permite
que evolucionen elásticamente con el campo de deformación local
\cite{Soneda2003}. Otra alternativa es promediar las ocupancias WS
sobre $N_f$ frames consecutivos:
\begin{equation}
  \bar \Omega_k \;=\; \frac{1}{N_f}\sum_{t} \Omega_k(t),
\end{equation}
y declarar vacancia los sitios con $\bar \Omega_k < 1/2$. Ambos
enfoques requieren más información (campo elástico o trayectoria
temporal) que el frame único usado por nuestro método.

\subsection{Posicionamiento del método propuesto}

\begin{table}[H]
\centering
\caption{Comparación cualitativa de los métodos discutidos.}
\label{tab:methods}
\small
\begin{tabular}{lccccc}
\toprule
Método & Ref.\ requerida & Robusto a $T$ & Cuenta vacancias & Costo & Detecta voids \\
\midrule
WS clásico                  & sí & no  & sí & $\mathcal{O}(N)$    & no       \\
WS adaptativo               & sí & medio & sí & $\mathcal{O}(N)$  & no       \\
Voronoi (OVITO)             & no & sí  & indirecto & $\mathcal{O}(N\log N)$ & sí \\
a-CNA \cite{Stukowski2012}  & no & sí  & no  & $\mathcal{O}(N)$  & parcial  \\
PTM \cite{Larsen2016}       & no & alto & no & $\mathcal{O}(N)$  & no       \\
Grid FaVaD \cite{Toussaint2020} & no & medio & parcial & $\mathcal{O}(N_\text{grid})$ & sí \\
SOAP fingerprint \cite{Toussaint2020} & sí & sí & no & $\mathcal{O}(N)$ & no \\
\textbf{Híbrido (este trabajo)} & opcional & alto & sí & $\mathcal{O}(N + N_\text{cand})$ & sí \\
\bottomrule
\end{tabular}
\end{table}

% =============================================================================
\section{Resumen y trabajo futuro}
% =============================================================================

El método híbrido propuesto satisface por construcción las cuatro
propiedades enunciadas en \S\ref{sec:ws-overest}: reproduce WS en
régimen de bajo daño (Proposición \ref{prop:ws-invariant}), es
robusto a ruido térmico (Proposición \ref{prop:thermal}), filtra
pares Frenkel ballísticos (Proposición \ref{prop:transit}) y admite
operación reference-free a través de las señales $s_\text{vor}$,
$s_\text{dens}$, $s_\text{topo}$.

\paragraph{Trabajo futuro identificado.}
\begin{enumerate}[label=(\arabic*)]
\item Reemplazar el proxy $s_\text{vor}$ por una integración real con
      \texttt{voro++} para capturar momentos de Voronoi de orden
      superior.
\item Incorporar $s_\text{topo}$ basado en PTM (con plantillas BCC,
      FCC, HCP, ico) para extender el método a aleaciones multifase.
\item Validación contra una base de ground-truth: cristales BCC y
      FCC con vacancias inyectadas en posiciones conocidas, MD a
      $T = \{0, 300, 600, 1000, 1500\}$~K, comparación cuantitativa
      $N_\text{vac}^\text{H}$ vs $N_\text{vac}^\text{WS}$ vs número
      inyectado.
\item Implementación de promedio temporal sobre múltiples frames
      para suprimir aún más la fluctuación instantánea.
\item Calibración Bayesiana de los pesos $w_\alpha$ mediante
      máxima verosimilitud sobre la base de ground-truth.
\end{enumerate}

% =============================================================================
\begin{thebibliography}{99}
% =============================================================================

\bibitem{Stukowski2010}
A.~Stukowski,
\textit{Structure identification methods for atomistic simulations of
crystalline materials},
Modelling and Simulation in Materials Science and Engineering
\textbf{20}, 045021 (2012).

\bibitem{Stukowski2010Ovito}
A.~Stukowski,
\textit{Visualization and analysis of atomistic simulation data with
OVITO – the Open Visualization Tool},
Modelling and Simulation in Materials Science and Engineering
\textbf{18}, 015012 (2010).

\bibitem{Stukowski2012}
A.~Stukowski,
\textit{Structure identification methods for atomistic simulations of
crystalline materials},
Modelling and Simulation in Materials Science and Engineering
\textbf{20}, 045021 (2012).

\bibitem{Larsen2016}
P.~M.~Larsen, S.~Schmidt, J.~Schiøtz,
\textit{Robust structural identification via polyhedral template matching},
Modelling and Simulation in Materials Science and Engineering
\textbf{24}, 055007 (2016).

\bibitem{Toussaint2020}
U.~von Toussaint, J.~Domínguez-Gutiérrez, M.~Compostella, M.~Rampp,
\textit{FaVaD: A software workflow for characterisation and visualizing
of defects in crystalline structures},
arXiv:2004.08184 (2020).

\bibitem{Nordlund2018}
K.~Nordlund et al.,
\textit{Primary radiation damage: A review of current understanding
and models},
Journal of Nuclear Materials \textbf{512}, 450 (2018).

\bibitem{Stoller2000}
R.~E.~Stoller,
\textit{The role of cascade energy and temperature in primary defect
formation in iron},
Journal of Nuclear Materials \textbf{276}, 22 (2000).

\bibitem{Soneda2003}
N.~Soneda, S.~Ishino, A.~Takahashi, K.~Dohi,
\textit{Modeling the microstructural evolution in bcc-Fe during
irradiation using kinetic Monte Carlo computer simulation},
Journal of Nuclear Materials \textbf{323}, 169 (2003).

\bibitem{Bjorkas2009}
C.~Björkas, K.~Nordlund,
\textit{Comparative study of cascade damage in Fe simulated with
recent potentials},
Nuclear Instruments and Methods B \textbf{259}, 853 (2009).

\bibitem{Becquart2010}
C.~S.~Becquart, C.~Domain,
\textit{Solute--point defect interactions in bcc systems: Focus on
first principles modelling in W and RPV steels},
Current Opinion in Solid State and Materials Science \textbf{16}, 115
(2012).

\bibitem{Granberg2016}
F.~Granberg et al.,
\textit{Mechanism of radiation damage reduction in equiatomic
multicomponent single-phase alloys},
Physical Review Letters \textbf{116}, 135504 (2016).

\bibitem{Norgett1975}
M.~J.~Norgett, M.~T.~Robinson, I.~M.~Torrens,
\textit{A proposed method of calculating displacement dose rates},
Nuclear Engineering and Design \textbf{33}, 50 (1975).

\bibitem{Brostow1978}
W.~Brostow, J.-P.~Dussault, B.~L.~Fox,
\textit{Construction of Voronoi polyhedra},
Journal of Computational Physics \textbf{29}, 81 (1978).

\bibitem{Rycroft2009}
C.~H.~Rycroft,
\textit{Voro++: A three-dimensional Voronoi cell library in C++},
Chaos \textbf{19}, 041111 (2009).

\bibitem{Silverman1986}
B.~W.~Silverman,
\textit{Density Estimation for Statistics and Data Analysis},
Chapman \& Hall (1986).

\bibitem{Kittel2004}
C.~Kittel,
\textit{Introduction to Solid State Physics},
8th ed., Wiley (2004).

\bibitem{Bartok2013}
A.~P.~Bartók, R.~Kondor, G.~Csányi,
\textit{On representing chemical environments},
Physical Review B \textbf{87}, 184115 (2013).

\bibitem{DominguezGutierrez2019}
J.~Domínguez-Gutiérrez, U.~von Toussaint,
\textit{On the detection and classification of material defects in
crystalline solids after energetic particle impact simulations},
Nuclear Materials and Energy \textbf{22}, 100724 (2019).

\end{thebibliography}

\end{document}
